% !PRINTED 206
\weberpage{206}
%<<CONT
betrachten, die dort mit $A''$ bezeichnet waren, und die, wenn man
$y_1 : y_2 = y$, $x_1 : x_2 = x$ setzt, wieder auf die Substitution (1)
führen. In (2) ist immer

\weq{3}{\alpha\delta - \beta\gamma = 1}

vorausgesetzt, aber die beiden Substitutionen
\[
\begin{pmatrix}\alpha, & \beta \\ \gamma, & \delta\end{pmatrix}, \quad
\begin{pmatrix}-\alpha, & -\beta \\ -\gamma, & -\delta\end{pmatrix}
\]
werden als identisch angesehen.

Es kann nicht zu einem Missverständniss führen, wenn wir
beide Arten von Substitutionen übereinstimmend durch ein
Symbol wie
\[
\Theta = \begin{pmatrix}\alpha, & \beta \\ \gamma, & \delta\end{pmatrix}
\]
bezeichnen, da ja dann in beiden die Compositionsregel genau
dieselbe ist.

Wenn nun

\weq{4}{a, a_1, a_2, \ldots, a_{\mu-1}}

ein System conjugirter Pole der Gruppe $G$ ist, so ist

\weq{5}{f(x) = (x-a)(x-a_1)\cdots(x-a_{\mu-1})}

eine ganze Function $\mu^{\text{ten}}$ Grades, deren Wurzeln jene conjugirten
Pole sind. Diese Function $f(x)$ hat folgende Eigenschaft:

Nehmen wir irgend eine Substitution der Gruppe $G$
\[
\psi(x) = \frac{\alpha x + \beta}{\gamma x + \delta},
\]
so sind die Wurzeln der Function

\weq{6}{(\gamma x + \delta)^{\mu} f[\psi(x)]}

die Grössen $\psi^{-1}(a), \psi^{-1}(a_1), \ldots, \psi^{-1}(a_{\mu-1})$. Diese aber stimmen,
von der Reihenfolge abgesehen, nach \S.\,52, 5. mit den
Grössen $a, a_1, a_2, \ldots, a_{\mu-1}$ überein, und folglich können sich die
Gleichungen (5) und (6) nur durch einen constanten Factor
unterscheiden. Wenn wir also

\weq{7}{x_2^{\mu} f\left(\frac{x_1}{x_2}\right) = f(x_1, x_2)}

setzen, so bleibt diese Form $f$, von einem constanten Factor abgesehen,
ungeändert, wenn auf $(x_1, x_2)$ irgend eine Substitution
der Gruppe $G$ angewandt wird.

Nach der Definition \S.\,40, 2. ist also $f(x_1, x_2)$ eine
\gesperrt{Invariante der Gruppe} $G$, und wir haben den Satz:
